%!TEX program = uplatex
\RequirePackage{plautopatch}
\documentclass[uplatex,dvipdfmx]{jlreq}
\usepackage{amsmath,amssymb}
\usepackage{enumerate}
\usepackage{url}

\pagestyle{empty}

\begin{document}

%======================================================================
% 講演1：国府寛司「双曲力学系」
%======================================================================
\begin{center}
{\large\bfseries 双曲力学系}\\[6pt]
\hspace*{\fill}{\large\bfseries 国府 寛司（京都大学 大学院理学研究科 数学教室）}
\end{center}
\vspace{1pc}

力学系とは決定論的な時間発展をするシステムを表す数学的概念であり、
常微分方程式（連続時間の力学系）や写像の反復合成（離散時間の力学系）
はその典型である。双曲力学系 (hyperbolic dynamical system) とは、
力学系が（少なくともその回帰的な部分において）伸びる方向と縮む方向に
「きれいに」分解されていることである。

力学系の双曲性の概念は、それ以前の先駆的な研究成果を踏まえて、
S. Smale によって一般的に定式化された。Smale は、構造安定な力学系、
すなわち摂動によって定性的構造を変えない力学系が力学系全体の中で
稠密であろうと考え、構造安定な力学系を双曲性によって特徴づけようと
した。この Smale のプログラムは60年代から70年代にかけて Smale と
その後継者たちによって盛んに研究され、構造安定な力学系の十分条件が
70年代半ばまでに得られた。また、それが（$C^1$-位相の下で）必要条件でも
あることが80年代に Mañé, Palis, 林によって証明された。
一方、構造安定な力学系が稠密でないことは70年代半ばまでには知られて
おり、また、気象学、生物学、電気工学などの自然科学や工学においては
カオスと呼ばれる複雑で予測困難な力学系の振舞いが発見されていたが、
これらの中には構造安定でないにもかかわらず摂動に対して安定に存在し
ているように見えるものがあり、構造不安定な力学系の研究はカオス
現象の理解の重要性は数学を超えて広く認識された。特に、構造安定性を
特徴づけるのは一様双曲性と呼ばれる最も厳格な双曲性であるが、構造
不安定な力学系においても、弱い形の双曲性（非一様な双曲性）は見られ、
それが力学系のカオスと呼ばれる複雑な振舞いに密接に関連していること
が次第にあきらかになってきた。

これに基づき、一般の構造不安定な力学系を何らかの理解可能な力学系で
近似するというアプローチが模索され、homoclinic 接触などの特徴的な
構造不安定性を持つ力学系が Palis--Takens、Newhouse らによって研究
された。今日ではそれを含む力学系のあるクラスが $C^1$-位相で稠密で
あろうという Palis の予想が提出され、このような立場からの力学系の
研究の指針となっている。

今回の Encounter with Mathematics では、以上のような力学系の
双曲性とそれに関連する話題をテーマに取り上げて、その基本的な部分を
概観したい。中心的な話題は
\begin{itemize}
\item 力学系の構造安定性が双曲性とどのように関連し特徴づけられるか？
\item 構造安定でない力学系においては何が起こりどのように理解されるか？
\end{itemize}
であり、前者は林氏の講演で、後者は浅岡氏の講演で解説される。また
このような双曲力学系や homoclinic 接触などを含む重要な例として
Hénon 写像と呼ばれる2次元の多項式写像の族があり、計算機による
数値実験も含めて多くの興味深い結果や予想が出され、一般的な研究の
試金石となっているが、三波氏の講演ではそれについての解説がなされる。

私の講演ではこれらの講演の前座として、力学系とその双曲性についての
最も基本的事項と、いくつかの重要な例を、主に離散時間の力学系の場合
について入門的な解説をする。具体的には
\begin{enumerate}
\item 力学系の基礎的概念と例
\item 力学系の双曲性（定義、安定／不安定多様体）
\item いくつかの重要な例（Morse--Smale 系、双曲的トーラス自己同型、馬蹄力学系）
\item 双曲力学系の基本的性質（スペクトル分解、マルコフ分割）
\end{enumerate}
などについて、アイディアを中心に述べるつもりである。

\bigskip
\noindent\textbf{参考文献}
\begin{enumerate}[{[1]}]
\item 国府　寛司 著；『力学系の基礎』、カオス全書2、朝倉書店、(2000)．
\item M. Shub；``Global Stability of Dynamical Systems'', Springer-Verlag, (1987).
（邦訳：『力学系』 上・下、シュプリンガー東京）
\item C. Robinson；``Dynamical Systems, Stability, Symbolic Dynamics, and Chaos'',
CRC Press, (1999).
\end{enumerate}

\noindent\textbf{主催者による註：}ここに文献[1]を挙げることを、
国府さんご自身はためらわれていたのですが、
主催者の知る限り、双曲力学系への入門書として類を見ないほどの
名著であると思いますので、
ご無礼を省みず挙げさせて頂きました。

\newpage

%======================================================================
% 講演2：林修平「力学系の安定性と通有性」
%======================================================================
\begin{center}
{\large\bfseries 力学系の安定性と通有性}\\[6pt]
\hspace*{\fill}{\large\bfseries 林 修平（東京大学 大学院数理科学研究科）}
\end{center}
\vspace{1pc}

双曲性理論における重要な概念として、通有性 (genericity) と安定性があります。
現代的な意味での最初の成果として、M. Peixoto による2次元の
フロー（連続時間の力学系）の研究がありますが、これは Morse--Smale 系
というフローが $C^1$ 位相の下で構造安定かつ稠密であるというものです。
つまり安定性は通有性でもあったわけです。Peixoto の結果の S. Smale に
よる高次元化の試みの中で Morse--Smale 系ではない馬蹄力学系が発見され、
理論は新たな局面を迎えますが、安定性の特徴づけへ向かって「公理 A」と
いう双曲性の概念が定式化されました。ところが、Peixoto の結果の高次
元化という Smale の夢は、安定でない力学系からなる開集合の発見により
崩壊します。以後、方向は安定性（構造安定性および $\Omega$-安定性）の
特徴づけの完成と、安定性に代わる新たな通有性を持つ有意義な性質の模索へ
と分岐することになります。安定性の方は、期待通りの成功をおさめ、
$C^1$ 位相については、離散時間と連続時間の両方の力学系で完全に解決され、
双曲性が安定性と本質的に同値であることが判明しました。この講演では、
解決までの流れを概観しながら、「$C^1$ 安定性予想」と呼ばれ長年未解決で
あった安定性の必要条件を中心に、証明のポイントを解説します。その中で
とりわけ重要な役割を果たしたのが Closing Lemma、Ergodic Closing Lemma、
Connecting Lemma と呼ばれる一連の $C^1$ 摂動定理ですが、これら $C^1$
摂動定理の証明のアイデアにも触れます。一方、通有性の方については、
非双曲性の領域にある力学系の持つ重要な現象としてホモクリニック接触と
heterodimensional cycle があり、J. Palis はこれらの現象が双曲性を妨げる
要因で、かつ、非双曲性の大局的理解への起点であると考えます。この立場
から提出された予想が、双曲性を持つ力学系にその2つの現象を持つ力学系を
合わせれば力学系全体の空間の中で稠密になるというものです。この Palis
予想への展望にも言及するつもりです。

\newpage

%======================================================================
% 講演3：浅岡正幸「ホモクリニック分岐」
%======================================================================
\begin{center}
{\large\bfseries ホモクリニック分岐}\\[6pt]
\hspace*{\fill}{\large\bfseries 浅岡 正幸（京都大学 大学院理学研究科 数学教室）}
\end{center}
\vspace{1pc}

J. Palis のプログラムに従えば、双曲力学系の世界の外側の大部分は
ホモクリニック分岐と呼ばれる分岐現象を通して理解できると思われる。
一方、ある力学系で起こる複雑な現象が、ホモクリニック分岐を通して
理解されるということも多い。
これらの理由から、ホモクリニック分岐を理解するために多くの研究が
なされてきたが、それらを通して明らかになってきたことの一つは、
まるで理解を拒絶するかのような非常に複雑なダイナミクスが
ホモクリニック分岐のそばでは豊富に現れるという、その野性的な側面である。
本講演では、ホモクリニック分岐の中でもよく調べられている
ホモクリニック接触について、そうした側面を中心に解説したい。

周期点の安定多様体と不安定多様体が横断的に交わらないとき、
その周期点はホモクリニック接触をもつという。
ホモクリニック接触は局所的な記述を持つため、その分岐の解析は
リスケーリングの方法によって局所的には2次写像や、
Hénon-like map の理論に還元することができる。
前半の講演では、これらホモクリニック接触を扱うための基本的な方法と、
そこから直接得られるいくつかの帰結について、また時間があれば、
ホモクリニック接触の開折における分岐集合の大きさを評価する
Palis--Takens 理論について解説したい。

S. Newhouse は1970年代に、リスケーリングと Cantor 集合の「厚み」の概念を
巧妙に用いることで、$r \geq 2$ について、
ホモクリニック接触を起こす $C^r$ 級力学系が一つあれば、
その任意の $C^r$-近傍に、ホモクリニック接触を持つような力学系を
稠密に持つような開集合（Newhouse 領域と呼ばれる）が存在することを示した。
後半の講演では、Newhouse 領域で generic に起こる野性的な
ダイナミクスに関する S. Newhouse や Y. Kaloshin の仕事の紹介をしたい。
時間が許せば、C. Bonatti--L. Díaz が発見した、Newhouse 現象とは
異なるメカニズムが引き起こす豊富な野性的ダイナミクスについての
結果や、ホモクリニック分岐がもたらす複雑さにも拘らず、
双曲力学系の外側の世界を理解しようとするいくつかの試みにも
触れたいと思う。

\bigskip
\noindent\textbf{参考文献：}
\begin{enumerate}[{[1]}]
\item J. Palis and J. Takens,
\textit{Hyperbolicity and sensitive chaotic dynamics at homoclinic bifurcations},
Cambridge Studies in Advanced Mathematics, 35.
Cambridge University Press, Cambridge, 1993.
\item C. Bonatti, L. Díaz, M. Viana,
\textit{Dynamics beyond uniform hyperbolicity},
Encyclopedia of Mathematical Sciences, 102.
Springer-Verlag, Berlin, 2004.
\end{enumerate}

\newpage

%======================================================================
% 講演4：三波篤郎「エノン写像」
%======================================================================
\begin{center}
{\large\bfseries エノン写像}\\[6pt]
\hspace*{\fill}{\large\bfseries 三波 篤郎（北見工業大学 情報システム工学科）}
\end{center}
\vspace{1pc}

エノン写像とは、次のような形をした $\mathbf{R}^2$ あるいは $\mathbf{C}^2$ から
それ自身への2次多項式で表わされる写像のことです。
\[H_{a,b}(x,y)=(-by+a-x^2,x)\]

2次多項式で表わされる $\mathbf{R}^2$ あるいは $\mathbf{C}^2$ からそれ自身への
写像で、逆写像がまた多項式であるもの（多項式自己同型写像）は、
アファイン写像による座標変換によって、
エノン写像かまたは、力学系的に単純なものに変換できることが知られています。
つまり、エノン写像は、最も単純な非線形の微分同相写像で、
力学系的に非自明なものの標準形なのです。

非線形系は、多くの場合、ホモクリニック接触というものを持っており、
その存在が、力学系的構造と分岐を極度に複雑なものにしています。
エノン写像はそれ自身が、いたるところに無数のホモクリニック接触を
持っていると同時に、ホモクリニック接触を局所的に見た時のモデルでもあります。
エノン写像は、ストレンジ・アトラクター（非双曲的アトラクター）を持つ
最も単純な写像であり、また $b=1$ の場合には、
面積保存写像の1パラメーター族となり、
KAM 理論的な構造と分岐を持つことが観察されます。

このように、エノン写像は、
力学系を考える上で、理解しなければならない基本的な対象の一つであると
言えますが、その数式の単純さとは裏腹に、
その構造や分岐は、未だに解明されてはいません。

この講演では、$\mathbf{R}^2$ 上のエノン写像の基本的な性質について説明し、
どの程度のことがこれまでに知られているかについて、
主に双曲性との関連という観点から解説します。

\bigskip
\noindent\textbf{参考文献：}
\begin{enumerate}[{[S]}]
\item[{[S]}] 三波篤郎、Hénon map について（第2版）、
(1998年度力学系勉強会 講演資料)\\
\url{http://www.math.sci.hokudai.ac.jp/~nami/dyn-ml/old-symp/1998toyama/}
\end{enumerate}

\end{document}
